% ============================================================
%  MediSync — Documentation Technique
%  Compilez avec : pdflatex documentation_technique.tex (x2)
% ============================================================
\documentclass[12pt,a4paper]{report}

% ---- Encodage & langue ----
\usepackage[utf8]{inputenc}
\usepackage[T1]{fontenc}
\usepackage[french]{babel}

% ---- Police ----
\usepackage{lmodern}
\usepackage{microtype}

% ---- Mise en page ----
\usepackage[top=2.5cm,bottom=2.5cm,left=2.8cm,right=2.8cm]{geometry}
\usepackage{setspace}
\onehalfspacing
\setlength{\parskip}{5pt}
\setlength{\parindent}{0pt}

% ---- Couleurs & graphiques ----
\usepackage[dvipsnames,svgnames,x11names]{xcolor}
\usepackage{graphicx}
\usepackage{tikz}
\usetikzlibrary{shapes.geometric,arrows.meta,positioning,calc,fit,backgrounds}

% ---- Tableaux ----
\usepackage{booktabs}
\usepackage{longtable}
\usepackage{tabularx}
\usepackage{multirow}
\usepackage{array}

% ---- Réduit l'espacement interne des tableaux ----
\setlength{\tabcolsep}{5pt}

% ---- Code & verbatim ----
\usepackage{listings}

% ---- En-têtes & pieds de page ----
\usepackage{fancyhdr}
\usepackage{emptypage}

% ---- Hyperliens ----
\usepackage[hidelinks,
  pdftitle={MediSync -- Documentation Technique},
  pdfauthor={Equipe MediSync},
  pdfsubject={Documentation technique complete}]{hyperref}

% ---- Autres utilitaires ----
\usepackage{enumitem}
\usepackage{tcolorbox}
\usepackage{caption}
\usepackage{amsmath}

\tcbuselibrary{skins,breakable}

% ============================================================
%  Palette de couleurs MediSync
% ============================================================
\definecolor{MedTeal}{HTML}{00838F}
\definecolor{MedTealDark}{HTML}{005F6B}
\definecolor{MedTealLight}{HTML}{E0F4F6}
\definecolor{DarkSlate}{HTML}{0F172A}
\definecolor{LightGray}{HTML}{F8FAFC}
\definecolor{BorderGray}{HTML}{CBD5E1}
\definecolor{CodeBg}{HTML}{F1F5F9}
\definecolor{SuccessGreen}{HTML}{166534}
\definecolor{WarningOrange}{HTML}{C2410C}
\definecolor{ErrorRed}{HTML}{DC2626}
\definecolor{CommentGray}{HTML}{64748B}
\definecolor{KeywordBlue}{HTML}{1D4ED8}

% ============================================================
%  Style des listings de code — tout en clair
% ============================================================
\lstdefinestyle{code}{
  basicstyle=\ttfamily\small\color{DarkSlate},
  backgroundcolor=\color{CodeBg},
  frame=single,
  rulecolor=\color{BorderGray},
  breaklines=true,
  breakatwhitespace=false,
  showstringspaces=false,
  keywordstyle=\color{MedTealDark}\bfseries,
  stringstyle=\color{SuccessGreen},
  commentstyle=\color{CommentGray}\itshape,
  columns=flexible,
  xleftmargin=10pt,
  xrightmargin=10pt,
  aboveskip=6pt,
  belowskip=6pt,
  tabsize=2,
}
\lstdefinestyle{bash}{
  style=code,
  language=bash,
  morekeywords={node,npm,ng,mongod,mkdir,cd,sudo,rm},
  morecomment=[l]{\#},
  keywordstyle=\color{MedTealDark}\bfseries,
}
\lstset{style=code}

% ============================================================
%  Boîtes personnalisées
% ============================================================
\newtcolorbox{infobox}[1][]{
  colback=MedTealLight,
  colframe=MedTeal,
  boxrule=0.8pt,
  arc=4pt,
  left=8pt, right=8pt, top=5pt, bottom=5pt,
  fonttitle=\bfseries\small\color{MedTealDark},
  coltitle=MedTealDark,
  #1
}
\newtcolorbox{warnbox}[1][]{
  colback=orange!8,
  colframe=WarningOrange,
  boxrule=0.8pt,
  arc=4pt,
  left=8pt, right=8pt, top=5pt, bottom=5pt,
  fonttitle=\bfseries\small\color{WarningOrange},
  coltitle=WarningOrange,
  #1
}

% ============================================================
%  En-têtes & pieds de page
% ============================================================
\pagestyle{fancy}
\fancyhf{}
\fancyhead[L]{\small\textcolor{MedTeal}{\textbf{MediSync}}%
  \textcolor{BorderGray}{\enspace|\enspace}%
  \textcolor{CommentGray}{Documentation Technique}}
\fancyhead[R]{\small\textcolor{CommentGray}{\nouppercase{\leftmark}}}
\fancyfoot[C]{\small\textcolor{CommentGray}{\thepage}}
\renewcommand{\headrulewidth}{0.4pt}
\renewcommand{\footrulewidth}{0pt}

% ============================================================
%  Commandes utilitaires
% ============================================================
\newcommand{\method}[1]{\texttt{\textcolor{MedTealDark}{\bfseries#1}}}
\newcommand{\role}[1]{%
  \texttt{\small\colorbox{MedTealLight}{\textcolor{MedTealDark}{#1}}}}
\newcommand{\field}[1]{\texttt{\small\textcolor{MedTealDark}{#1}}}
\newcommand{\typ}[1]{\textit{\small\textcolor{CommentGray}{#1}}}

% ============================================================
%  DOCUMENT
% ============================================================
\begin{document}

% ============================================================
%  PAGE DE GARDE — fond blanc, accents teal
% ============================================================
\begin{titlepage}
  \pagecolor{white}
  \color{DarkSlate}

  \vspace*{1.5cm}

  % Bande décorative supérieure
  \noindent\textcolor{MedTeal}{\rule{\textwidth}{4pt}}

  \vspace{0.5cm}

  \begin{center}
    % Logo circulaire
    \begin{tikzpicture}
      \fill[MedTeal] (0,0) circle (1.6cm);
      \fill[MedTealLight] (0,0) circle (1.3cm);
      \node[font=\fontsize{28}{28}\selectfont\bfseries\color{MedTealDark}]
        at (0,0) {M};
    \end{tikzpicture}

    \vspace{0.7cm}
    {\fontsize{36}{42}\selectfont\bfseries\textcolor{MedTeal}{MediSync}}

    \vspace{0.2cm}
    {\large\textcolor{CommentGray}{Système de Gestion de Clinique Médicale}}

    \vspace{1.2cm}

    % Boîte titre centrale
    \noindent\colorbox{MedTealLight}{%
      \parbox{0.85\textwidth}{%
        \vspace{0.5cm}
        \centering
        {\LARGE\bfseries\textcolor{MedTealDark}{Documentation Technique Complète}}
        \vspace{0.5cm}
      }%
    }

    \vspace{1.5cm}

    % Table d'informations
    \begin{tabular}{@{} r @{\quad} l @{}}
      \textcolor{CommentGray}{Version} &
        \textcolor{DarkSlate}{\textbf{1.0.0}} \\[6pt]
      \textcolor{CommentGray}{Date} &
        \textcolor{DarkSlate}{\textbf{Mai 2026}} \\[6pt]
      \textcolor{CommentGray}{Frontend} &
        \textcolor{DarkSlate}{\textbf{Angular 17 + Angular Material}} \\[6pt]
      \textcolor{CommentGray}{Backend} &
        \textcolor{DarkSlate}{\textbf{Node.js 18 + Express 4}} \\[6pt]
      \textcolor{CommentGray}{Base de données} &
        \textcolor{DarkSlate}{\textbf{MongoDB 7 (Mongoose 8)}} \\[6pt]
    \end{tabular}
  \end{center}

  \vfill

  \noindent\textcolor{BorderGray}{\rule{\textwidth}{1pt}}
  \vspace{0.2cm}
  \begin{center}
    \small\textcolor{CommentGray}{Projet académique — Soutenance juin 2026}
  \end{center}
  \vspace{0.3cm}
  \noindent\textcolor{MedTeal}{\rule{\textwidth}{4pt}}
\end{titlepage}
\nopagecolor

% ============================================================
%  TABLE DES MATIÈRES
% ============================================================
\tableofcontents
\thispagestyle{empty}
\clearpage

% ============================================================
%  CHAPITRE 1 — PRÉSENTATION DU PROJET
% ============================================================
\chapter{Présentation du Projet}

\section{Contexte et Objectifs}

\textbf{MediSync} est une application web full-stack de gestion de clinique médicale développée dans le cadre d'un projet académique. Elle vise à digitaliser et centraliser l'ensemble des flux d'une structure de soins : prise de rendez-vous, suivi des dossiers médicaux, facturation, gestion du personnel et pilotage analytique.

L'application répond aux exigences du cahier des charges en supportant \textbf{quatre rôles utilisateurs distincts} avec des interfaces et permissions adaptées à chaque profil professionnel.

\section{Rôles Utilisateurs}

\begin{center}
\small
\begin{tabular}{@{} p{2.6cm} p{11.5cm} @{}}
  \toprule
  \textbf{Rôle} & \textbf{Responsabilités principales} \\
  \midrule
  \role{patient} &
    Tableau de bord, recherche et réservation de médecins, suivi des rendez-vous,
    dossier médical personnel, factures, avis post-consultation. \\[4pt]
  \role{medecin} &
    Planning journalier, saisie des comptes rendus (ordonnances, prescriptions,
    pièces jointes), gestion des congés et indisponibilités, profil. \\[4pt]
  \role{secretaire} &
    Planning global (confirmation, no-show, report), création de dossiers patients,
    émission et suivi des factures. \\[4pt]
  \role{administrateur} &
    Tableau de bord KPI, CRUD personnel, configuration clinique, export analytique
    (PDF/XLSX), journal d'audit, 2FA. \\
  \bottomrule
\end{tabular}
\end{center}

\section{Fonctionnalités Clés}

\begin{itemize}[leftmargin=*,itemsep=2pt]
  \item \textbf{Authentification multi-modale} : identifiant/mot de passe, Google OAuth 2.0, TOTP 2FA (admin).
  \item \textbf{Gestion des rendez-vous} : créneaux en temps réel, six statuts métier.
  \item \textbf{Dossier médical électronique} : antécédents, allergies, consultations, ordonnances, pièces jointes (PDF, JPG, PNG, DICOM $\leq$ 20 Mo).
  \item \textbf{Facturation} : émission, feuille de soins (Sécu), PDF, envoi e-mail, impayés.
  \item \textbf{Notifications} : rappels automatiques de rendez-vous (cron horaire, Resend).
  \item \textbf{Analytics} : KPIs, revenus, exports PDF/XLSX (PDFKit, ExcelJS).
  \item \textbf{Journal d'audit} : traçabilité de chaque accès sensible (RGPD).
\end{itemize}

% ============================================================
%  CHAPITRE 2 — ARCHITECTURE
% ============================================================
\chapter{Architecture de l'Application}

\section{Vue d'Ensemble}

MediSync adopte une architecture \textbf{découplée client-serveur} (SPA + REST API) :

\begin{itemize}[leftmargin=*,itemsep=2pt]
  \item \textbf{Frontend SPA} (Angular 17) — communique avec le backend via requêtes HTTP authentifiées (JWT Bearer).
  \item \textbf{Backend REST API} (Express 4 / Node 18) — logique métier et accès aux données.
  \item \textbf{Base de données} (MongoDB 7) — documents JSON ; Mongoose 8 fournit l'ODM, la validation et les hooks.
  \item \textbf{Stockage fichiers} — dossier \texttt{uploads/} servi en statique par Express.
  \item \textbf{Service e-mail} — API Resend (HTTPS) pour les notifications transactionnelles.
\end{itemize}

\vspace{0.4cm}

% ---- Diagramme architecture : disposition verticale + services latéraux ----
\begin{center}
\begin{tikzpicture}[
  box/.style={
    draw=MedTeal, fill=MedTealLight, rounded corners=5pt,
    minimum width=3.6cm, minimum height=0.85cm,
    align=center, font=\small\bfseries\color{MedTealDark}
  },
  extbox/.style={
    draw=BorderGray, fill=LightGray, rounded corners=5pt,
    minimum width=3.2cm, minimum height=0.75cm,
    align=center, font=\small\color{CommentGray}
  },
  arr/.style={-Latex, thick, MedTeal},
  darr/.style={-Latex, thick, dashed, BorderGray},
  lbl/.style={font=\footnotesize\itshape, text=CommentGray, fill=white,
              inner sep=1pt}
]
  % Colonne principale (verticale, centrée)
  \node[box] (browser) at (0,0)       {Navigateur Web};
  \node[box] (ng)      at (0,-1.8)    {Angular 17\\{\footnotesize Port 4200}};
  \node[box] (api)     at (0,-3.6)    {Express API\\{\footnotesize Port 3000}};
  \node[box] (db)      at (0,-5.4)    {MongoDB\\{\footnotesize Port 27017}};

  % Flèches bidirectionnelles colonne principale
  \draw[arr] (browser) -- node[lbl,right]{HTTP / SPA} (ng);
  \draw[arr] (ng)      -- node[lbl,right]{REST / JWT} (api);
  \draw[arr] (api)     -- node[lbl,right]{Mongoose ODM} (db);
  \draw[darr,<->] (db.north) -- +(0,0.4)
    node[lbl,left,pos=1]{Documents JSON} -- (api.south);

  % Service externe : Google OAuth (à gauche)
  \node[extbox] (google) at (-5.2,-1.8) {Google OAuth 2.0};
  \draw[darr] (ng.west) -- node[lbl,above]{\footnotesize OAuth} (google.east);

  % Service externe : Resend (à droite)
  \node[extbox] (resend) at (5.2,-3.6) {Resend\\{\footnotesize E-mail API}};
  \draw[darr] (api.east) -- node[lbl,above]{\footnotesize HTTPS} (resend.west);

  % Étiquettes de section
  \node[font=\footnotesize\bfseries\color{CommentGray}, above=0.15cm of browser]
    {Couche Présentation};
  \node[font=\footnotesize\bfseries\color{CommentGray}, below=0.15cm of db]
    {Couche Persistance};
\end{tikzpicture}
\end{center}

\section{Structure des Répertoires}

\subsection{Backend (\texttt{/backend})}

\begin{lstlisting}[style=bash]
backend/
  server.js             # Entree : middlewares globaux, montage des routes
  seed.js               # Donnees de demonstration
  src/
    config/db.js        # Connexion Mongoose
    controllers/        # Logique metier (9 controleurs)
    middleware/
      auth.js           # JWT verify -> req.user
      role.js           # authorize(...roles) -> 403 si role absent
      upload.js         # multer : filtre MIME, limite 20 Mo
      validate.js       # express-validator : arrete si erreurs
    models/             # 9 schemas Mongoose
    routes/             # 9 routeurs Express
    utils/
      auditLogger.js    # Ecrit dans AuditLog (actions sensibles)
      emailService.js   # Resend : sendNotification()
      emailTemplates.js # Templates HTML des e-mails
      generateToken.js  # jwt.sign() avec expiry
      pdfGenerator.js   # PDFKit : factures, ordonnances, feuilles de soins
      reminderScheduler.js  # node-cron horaire
  uploads/              # Pieces jointes medicales (statique)
\end{lstlisting}

\subsection{Frontend (\texttt{/frontend/src})}

\begin{lstlisting}[style=bash]
frontend/src/app/
  app.config.ts         # provideRouter, HttpClient, locale FR, theme
  app.routes.ts         # Routes racines avec lazy loading
  core/
    guards/             # auth.guard.ts, role.guard.ts
    interceptors/       # auth.interceptor.ts, error.interceptor.ts
    models/index.ts     # Interfaces TypeScript (Account, Appointment...)
    services/           # 10 services Angular (wrappers HTTP)
    utils/              # token-storage.ts, date.ts
    data/               # medications.ts, consultation-templates.ts
  features/
    auth/               # login, register, 2FA, complete-profile
    patient/            # dashboard, search-doctors, book-appointment,
                        # my-appointments, my-record, my-invoices
    practitioner/       # calendar, consultation-form, leaves, profile
    secretary/          # schedule, patients, invoices + dialogs
    admin/              # dashboard, staff, facility, analytics,
                        # audit-log, 2FA-setup
\end{lstlisting}

\section{Flux d'Authentification}

\begin{center}
\begin{tikzpicture}[
  node distance=0.7cm,
  step/.style={
    draw=MedTeal, rounded corners=3pt, fill=MedTealLight,
    minimum width=8.5cm, minimum height=0.65cm,
    align=center, font=\small\color{DarkSlate}
  },
  dec/.style={
    draw=WarningOrange, diamond, fill=orange!10,
    minimum width=1.6cm, minimum height=0.65cm,
    align=center, font=\small\bfseries\color{WarningOrange}
  },
  altbox/.style={
    draw=WarningOrange, rounded corners=3pt, fill=orange!10,
    minimum width=4cm, minimum height=0.65cm,
    align=center, font=\small\color{DarkSlate}
  },
  arr/.style={-Latex, thick, MedTeal},
  warr/.style={-Latex, thick, dashed, WarningOrange},
  lbl/.style={font=\footnotesize\color{CommentGray}, fill=white, inner sep=1pt}
]
  \node[step]  (login)    {1.\ POST \texttt{/api/auth/login} (identifier + password)};
  \node[step, below=of login]  (bcrypt)  {2.\ Vérification \texttt{bcrypt} du mot de passe};
  \node[dec,  below=0.6cm of bcrypt]  (has2fa)  {2FA?};
  \node[step, below=1.3cm of has2fa]  (jwt)     {3a.\ Génération JWT (7j) + payload \texttt{\{id, role\}}};
  \node[step, below=of jwt]    (ls)      {4.\ Stockage \texttt{localStorage} (token + user)};
  \node[step, below=of ls]     (guard)   {5.\ AuthGuard / RoleGuard $\rightarrow$ redirection portail};
  \node[step, below=of guard]  (icp)     {6.\ AuthInterceptor : \texttt{Authorization: Bearer <token>}};

  % Branche 2FA (à droite du diamant)
  \node[altbox, right=1.2cm of has2fa] (totp)
    {3b.\ \texttt{tempToken} émis\\POST \texttt{/api/auth/2fa/verify}};

  % Flèches principales
  \draw[arr] (login)  -- (bcrypt);
  \draw[arr] (bcrypt) -- (has2fa);
  \draw[arr] (has2fa) -- node[lbl,left]{Non} (jwt);
  \draw[arr] (jwt)    -- (ls);
  \draw[arr] (ls)     -- (guard);
  \draw[arr] (guard)  -- (icp);

  % Branche 2FA
  \draw[warr] (has2fa) -- node[lbl,above]{Oui} (totp);
  \draw[warr] (totp.south) |- (jwt.east);
\end{tikzpicture}
\end{center}

% ============================================================
%  CHAPITRE 3 — PILE TECHNOLOGIQUE
% ============================================================
\chapter{Pile Technologique}

\section{Frontend}

\begin{center}
\small
\begin{tabular}{@{} l l l p{6cm} @{}}
  \toprule
  \textbf{Technologie} & \textbf{Version} & \textbf{Rôle} & \textbf{Justification} \\
  \midrule
  Angular              & 17.3  & Framework SPA     & Standalone components, routing, DI \\
  Angular Material     & 17.3  & Bibliothèque UI   & Composants accessibles, thème Material 3 \\
  Angular CDK          & 17.3  & Utilitaires UI    & Overlay, virtual scroll \\
  TypeScript           & 5.4   & Langage           & Typage statique, interfaces métier \\
  RxJS                 & 7.8   & Réactivité        & Observables pour les appels HTTP \\
  Chart.js / ng2-charts& 4.5/5.0 & Graphiques      & KPIs et analytics admin \\
  SCSS                 & —     & Styles            & Variables, thème teal personnalisé \\
  \bottomrule
\end{tabular}
\end{center}

\section{Backend}

\begin{center}
\small
\begin{tabular}{@{} l l l p{5.5cm} @{}}
  \toprule
  \textbf{Bibliothèque}   & \textbf{Version} & \textbf{Rôle}         & \textbf{Description} \\
  \midrule
  Express                 & 4.19  & Framework HTTP     & Routing, middlewares, serveur REST \\
  Mongoose                & 8.4   & ODM MongoDB        & Schémas, validation, hooks pre/post \\
  bcryptjs                & 2.4   & Sécurité           & Hachage des mots de passe (salt 10) \\
  jsonwebtoken            & 9.0   & Auth               & Génération / vérification JWT HS256 \\
  express-validator       & 7.2   & Validation         & Validation des corps de requêtes \\
  speakeasy               & 2.0   & 2FA TOTP           & Secret + vérification code 6 chiffres \\
  qrcode                  & 1.5   & 2FA                & QR Code pour apps authenticator \\
  google-auth-library     & 9.11  & OAuth              & Vérification \texttt{idToken} Google \\
  multer                  & 2.1   & Upload             & Stockage local, filtre MIME, 20 Mo max \\
  PDFKit                  & 0.18  & PDF                & Factures, ordonnances, feuilles de soins \\
  ExcelJS                 & 4.4   & Export             & Feuilles de calcul XLSX \\
  Resend                  & 6.12  & E-mail             & Notifications transactionnelles \\
  node-cron               & 4.2   & Tâches planifiées  & Rappels RDV (exécution horaire) \\
  morgan                  & 1.10  & Logging            & Journal des requêtes HTTP \\
  cors                    & 2.8   & Sécurité           & Restriction CORS à l'URL du frontend \\
  dotenv                  & 16.4  & Configuration      & Variables d'environnement \\
  \bottomrule
\end{tabular}
\end{center}

\section{Infrastructure}

\begin{center}
\small
\begin{tabular}{@{} l l l @{}}
  \toprule
  \textbf{Composant}        & \textbf{Technologie}  & \textbf{Version minimale} \\
  \midrule
  Runtime serveur           & Node.js               & 18 LTS \\
  Base de données           & MongoDB               & 6+ \\
  Gestionnaire de paquets   & npm                   & 9+ \\
  CLI Angular               & @angular/cli          & 17.3 \\
  \bottomrule
\end{tabular}
\end{center}

% ============================================================
%  CHAPITRE 4 — SCHÉMA DE BASE DE DONNÉES
% ============================================================
\chapter{Schéma de Base de Données}

MediSync utilise \textbf{MongoDB} (NoSQL orienté documents). Les schémas Mongoose définissent \textbf{neuf collections}. Les relations inter-documents sont matérialisées par des références \texttt{ObjectId} populées côté backend.

\section{Vue Relationnelle des Collections}

\begin{center}
\begin{tikzpicture}[
  coll/.style={
    draw=MedTeal, fill=MedTealLight, rounded corners=4pt,
    minimum width=2.8cm, minimum height=0.75cm,
    align=center, font=\small\bfseries\color{MedTealDark}
  },
  ref/.style={-Latex, dashed, CommentGray, font=\scriptsize,
              text=CommentGray},
]
  % Ligne 1 : pivot Account centré
  \node[coll] (acc) at (0,0) {Account};

  % Ligne 2 : profils
  \node[coll] (dp) at (-4,-2)  {DoctorProfile};
  \node[coll] (pp) at (0,-2)   {PatientProfile};
  \node[coll] (sp) at (4,-2)   {SecretaryProfile};

  % Ligne 3 : entités transactionnelles
  \node[coll] (apt) at (-2.5,-4) {Appointment};
  \node[coll] (mr)  at (2.5,-4)  {MedicalRecord};

  % Ligne 4 : entités dérivées
  \node[coll] (inv) at (-4.5,-6) {Invoice};
  \node[coll] (rv)  at (0,-6)    {Review};
  \node[coll] (al)  at (4.5,-6)  {AuditLog};

  % Relations vers Account
  \draw[ref] (dp) -- node[left]{\scriptsize account} (acc);
  \draw[ref] (pp) -- node[right,near start]{\scriptsize account} (acc);
  \draw[ref] (sp) -- node[right]{\scriptsize account} (acc);

  % Relations vers profils
  \draw[ref] (apt) -- node[left]{\scriptsize doctor} (dp);
  \draw[ref] (apt) -- node[right,near start]{\scriptsize patient} (pp);
  \draw[ref] (mr)  -- node[right]{\scriptsize patient} (pp);

  % Relations vers Appointment / profils
  \draw[ref] (inv) -- node[left]{\scriptsize appointment} (apt);
  \draw[ref] (inv.north east) -- node[above,sloped]{\scriptsize patient} (pp.south);
  \draw[ref] (rv)  -- node[right,near start]{\scriptsize appointment} (apt);

  % AuditLog → Account
  \draw[ref] (al) -- node[right]{\scriptsize userAccount} (acc);
\end{tikzpicture}
\end{center}

\section{Descriptions Détaillées des Collections}

% ---- Account ----
\subsection{Account}

Pivot central de l'identité — représente le compte d'authentification de tout utilisateur.

\begin{center}
\small
\begin{tabularx}{\textwidth}{@{} l l X @{}}
  \toprule
  \textbf{Champ} & \textbf{Type} & \textbf{Description} \\
  \midrule
  \field{email}               & \typ{String}  & E-mail (unique, sparse, lowercase). \\
  \field{socialSecurityNumber}& \typ{String}  & Numéro SS (unique, sparse). \\
  \field{password}            & \typ{String}  & Haché bcrypt (salt 10). Requis si pas \texttt{googleId}. \\
  \field{googleId}            & \typ{String}  & Identifiant Google OAuth 2.0. \\
  \field{role}                & \typ{Enum}    & \texttt{patient | medecin | secretaire | administrateur} \\
  \field{twoFactorSecret}     & \typ{String}  & Secret TOTP speakeasy (admins uniquement). \\
  \field{profileCompleted}    & \typ{Boolean} & \texttt{false} pour les patients OAuth en attente de complétion. \\
  \field{createdAt / updatedAt}& \typ{Date}   & Timestamps automatiques Mongoose. \\
  \bottomrule
\end{tabularx}
\end{center}

\begin{infobox}[title=Contrainte d'identifiant]
  Un compte doit posséder au moins un identifiant parmi : \texttt{email}, \texttt{socialSecurityNumber} ou \texttt{googleId}. Vérifié dans le hook \texttt{pre('validate')}.
\end{infobox}

% ---- DoctorProfile ----
\subsection{DoctorProfile}

Profil enrichi du médecin, lié à son \texttt{Account} via \texttt{account: ObjectId}.

\begin{center}
\small
\begin{tabularx}{\textwidth}{@{} l l X @{}}
  \toprule
  \textbf{Champ} & \textbf{Type} & \textbf{Description} \\
  \midrule
  \field{account}             & \typ{ObjectId$\to$Account}    & Compte d'authentification. \\
  \field{firstName / lastName}& \typ{String}                  & Identité civile. \\
  \field{specialties}         & \typ{[String]}                & Ex. : \texttt{["Cardiologie"]}. \\
  \field{languages}           & \typ{[String]}                & Langues parlées. \\
  \field{location}            & \typ{String}                  & Localisation dans l'établissement. \\
  \field{sector}              & \typ{Number (1|2|3)}          & Secteur tarifaire CPAM. \\
  \field{baseFee}             & \typ{Number}                  & Tarif de base de la consultation (€). \\
  \field{fees}                & \typ{[\{code, label, price\}]}& Actes spécifiques (CCAM). \\
  \field{schedule.workDays}   & \typ{[Enum]}                  & Jours de travail. \\
  \field{schedule.startHour / endHour} & \typ{String HH:mm}  & Plage horaire par défaut. \\
  \field{schedule.defaultDuration}& \typ{Number (15|30|60)}   & Durée standard en minutes. \\
  \field{leaves}              & \typ{[Sous-doc]}              & Congés : startDate, endDate, reason. \\
  \bottomrule
\end{tabularx}
\end{center}

% ---- PatientProfile ----
\subsection{PatientProfile}

\begin{center}
\small
\begin{tabularx}{\textwidth}{@{} l l X @{}}
  \toprule
  \textbf{Champ} & \textbf{Type} & \textbf{Description} \\
  \midrule
  \field{account}             & \typ{ObjectId$\to$Account}  & Compte d'authentification. \\
  \field{firstName / lastName}& \typ{String}                & Identité civile. \\
  \field{dateOfBirth}         & \typ{Date}                  & Date de naissance. \\
  \field{phoneNumber}         & \typ{String}                & Numéro de téléphone. \\
  \field{dependents}          & \typ{[Sous-doc]}            & Ayants droit (enfant, conjoint, parent). \\
  \field{dependents[].relation}& \typ{Enum}                 & \texttt{enfant | conjoint | parent} \\
  \field{dependents[].allergies}& \typ{[String]}            & Allergies propres à l'ayant droit. \\
  \bottomrule
\end{tabularx}
\end{center}

% ---- Appointment ----
\subsection{Appointment}

\begin{center}
\small
\begin{tabularx}{\textwidth}{@{} l l X @{}}
  \toprule
  \textbf{Champ} & \textbf{Type} & \textbf{Description} \\
  \midrule
  \field{patient}             & \typ{ObjectId$\to$PatientProfile} & Null pour les blocs \texttt{indisponible}. \\
  \field{doctor}              & \typ{ObjectId$\to$DoctorProfile}  & Médecin concerné. \\
  \field{dependentId}         & \typ{ObjectId}                    & Sous-doc ayant droit, si applicable. \\
  \field{startTime / endTime} & \typ{Date}                        & Plage temporelle du rendez-vous. \\
  \field{duration}            & \typ{Number (15|30|60)}           & Durée en minutes. \\
  \field{reason}              & \typ{String}                      & Motif de consultation. \\
  \field{notes}               & \typ{String}                      & Remarques libres (secrétaire). \\
  \field{room}                & \typ{String}                      & Salle affectée (\texttt{Facility.rooms[].roomName}). \\
  \field{status}              & \typ{Enum}                        &
    \texttt{en attente | confirmé | annulé | terminé |}
    \texttt{indisponible | no-show} \\
  \bottomrule
\end{tabularx}
\end{center}

Index composites : \texttt{\{doctor, startTime\}} et \texttt{\{patient, startTime desc\}}.

% ---- MedicalRecord ----
\subsection{MedicalRecord}

Un seul document par patient (contrainte \texttt{unique} sur \texttt{patient}).

\begin{center}
\small
\begin{tabularx}{\textwidth}{@{} l l X @{}}
  \toprule
  \textbf{Champ} & \textbf{Type} & \textbf{Description} \\
  \midrule
  \field{patient}                     & \typ{ObjectId$\to$PatientProfile}  & Référence unique. \\
  \field{history}                     & \typ{[String]}                     & Antécédents médicaux. \\
  \field{allergies}                   & \typ{[String]}                     & Allergies connues. \\
  \field{consultations[].appointmentId}& \typ{ObjectId$\to$Appointment}    & RDV source. \\
  \field{consultations[].doctorId}    & \typ{ObjectId$\to$DoctorProfile}   & Médecin rédacteur. \\
  \field{consultations[].report}      & \typ{String}                       & Compte rendu structuré. \\
  \field{consultations[].prescriptions}& \typ{[Sous-doc]}                  & Médicament, posologie, durée. \\
  \field{attachments[].fileUrl}       & \typ{String}                       & Chemin serveur du fichier. \\
  \field{attachments[].fileType}      & \typ{Enum}                         & \texttt{PDF | JPG | PNG | DICOM} \\
  \field{attachments[].uploadedBy}    & \typ{ObjectId$\to$Account}         & Auteur du dépôt. \\
  \bottomrule
\end{tabularx}
\end{center}

% ---- Invoice ----
\subsection{Invoice}

\begin{center}
\small
\begin{tabularx}{\textwidth}{@{} l l X @{}}
  \toprule
  \textbf{Champ} & \textbf{Type} & \textbf{Description} \\
  \midrule
  \field{appointment} & \typ{ObjectId$\to$Appointment}   & RDV facturé. \\
  \field{patient}     & \typ{ObjectId$\to$PatientProfile} & Patient facturé. \\
  \field{doctor}      & \typ{ObjectId$\to$DoctorProfile}  & Médecin praticien. \\
  \field{amount}      & \typ{Number}                      & Montant en euros. \\
  \field{nomenclature}& \typ{String}                      & Code acte (ex. : \textit{Secteur 1 – CCAM}). \\
  \field{status}      & \typ{Enum}                        & \texttt{payé | en attente | impayé} \\
  \field{issuedAt}    & \typ{Date}                        & Date d'émission (défaut : maintenant). \\
  \field{paidAt}      & \typ{Date}                        & Date de paiement effectif. \\
  \bottomrule
\end{tabularx}
\end{center}

Index : \texttt{\{patient, status\}}.

% ---- Facility ----
\subsection{Facility}

Collection singleton — un seul document représente la clinique.

\begin{center}
\small
\begin{tabularx}{\textwidth}{@{} l l X @{}}
  \toprule
  \textbf{Champ} & \textbf{Type} & \textbf{Description} \\
  \midrule
  \field{name}                    & \typ{String}   & Nom de la clinique. \\
  \field{address}                 & \typ{String}   & Adresse postale complète. \\
  \field{contactEmail / Phone}    & \typ{String}   & Coordonnées de l'établissement. \\
  \field{specialtiesOffered}      & \typ{[String]} & Spécialités disponibles. \\
  \field{openingHours}            & \typ{String}   & Plage d'ouverture (texte libre). \\
  \field{rooms[].roomName}        & \typ{String}   & Nom de la salle (ex. : \textit{Radiologie 1}). \\
  \field{rooms[].equipment}       & \typ{[String]} & Équipements associés à la salle. \\
  \bottomrule
\end{tabularx}
\end{center}

% ---- AuditLog ----
\subsection{AuditLog}

\begin{center}
\small
\begin{tabularx}{\textwidth}{@{} l l X @{}}
  \toprule
  \textbf{Champ} & \textbf{Type} & \textbf{Description} \\
  \midrule
  \field{userAccount} & \typ{ObjectId$\to$Account} & Auteur de l'action. \\
  \field{action}      & \typ{Enum}                 &
    \texttt{CONNEXION}, \texttt{ACCES\_DOSSIER},
    \texttt{MODIFICATION\_DOSSIER},
    \texttt{CREATION\_COMPTE},
    \texttt{EXPORT\_DONNEES} \\
  \field{targetId}    & \typ{String}               & ID de la ressource cible. \\
  \field{ipAddress}   & \typ{String}               & IP de la requête. \\
  \field{details}     & \typ{String}               & Description libre de l'action. \\
  \bottomrule
\end{tabularx}
\end{center}

% ---- Review ----
\subsection{Review}

\begin{center}
\small
\begin{tabularx}{\textwidth}{@{} l l X @{}}
  \toprule
  \textbf{Champ} & \textbf{Type} & \textbf{Description} \\
  \midrule
  \field{appointment}  & \typ{ObjectId$\to$Appointment}   & RDV évalué (unique). \\
  \field{patient}      & \typ{ObjectId$\to$PatientProfile} & Auteur de l'avis. \\
  \field{doctor}       & \typ{ObjectId$\to$DoctorProfile}  & Médecin évalué. \\
  \field{rating}       & \typ{Number (1--5)}               & Note de satisfaction. \\
  \field{comment}      & \typ{String}                      & Commentaire libre. \\
  \field{isIssueReport}& \typ{Boolean}                     & Signalement grave (défaut : false). \\
  \bottomrule
\end{tabularx}
\end{center}

% ============================================================
%  CHAPITRE 5 — API REST
% ============================================================
\chapter{Référence API REST}

Toutes les routes sont préfixées par \texttt{/api}. La base URL par défaut est \texttt{http://localhost:3000}.

\begin{infobox}[title=Convention d'authentification]
  Toutes les routes protégées requièrent l'en-tête :
  \texttt{Authorization: Bearer <jwt\_token>}\\
  Un token absent ou expiré retourne \texttt{401 Unauthorized}.
\end{infobox}

\section{Authentification — \texttt{/api/auth}}

\begin{center}
\small
\begin{tabular}{@{} l p{5.2cm} p{2.8cm} p{4.5cm} @{}}
  \toprule
  \textbf{Méth.} & \textbf{Endpoint} & \textbf{Accès} & \textbf{Description} \\
  \midrule
  \method{POST}  & \texttt{/register}         & Public       & Crée un compte (role, email/SSN, password). \\
  \method{POST}  & \texttt{/login}            & Public       & Connexion par identifiant + password. Retourne JWT ou \texttt{tempToken} si 2FA activée. \\
  \method{POST}  & \texttt{/google}           & Public       & Connexion Google OAuth (\texttt{idToken}). \\
  \method{POST}  & \texttt{/2fa/verify}       & Public       & Vérifie le code TOTP, retourne le JWT final. \\
  \method{GET}   & \texttt{/me}               & Authentifié  & Profil du compte connecté. \\
  \method{PATCH} & \texttt{/me}               & Authentifié  & Mise à jour du profil patient. \\
  \method{POST}  & \texttt{/complete-profile} & Authentifié  & Complétion profil (patients Google OAuth). \\
  \method{POST}  & \texttt{/2fa/setup}        & Admin        & Génère secret TOTP + QR Code. \\
  \method{POST}  & \texttt{/2fa/confirm-setup}& Admin        & Active la 2FA après vérification du premier code. \\
  \method{POST}  & \texttt{/2fa/disable}      & Admin        & Désactive la 2FA. \\
  \bottomrule
\end{tabular}
\end{center}

\section{Rendez-vous — \texttt{/api/appointments}}

\begin{center}
\small
\begin{tabular}{@{} l p{4.5cm} p{3.8cm} p{4.5cm} @{}}
  \toprule
  \textbf{Méth.} & \textbf{Endpoint} & \textbf{Rôles} & \textbf{Description} \\
  \midrule
  \method{POST}  & \texttt{/}                  & patient, secretaire, admin & Crée un rendez-vous. \\
  \method{GET}   & \texttt{/}                  & secretaire, admin          & Planning complet. \\
  \method{GET}   & \texttt{/my-appointments}   & patient                    & Historique du patient connecté. \\
  \method{GET}   & \texttt{/available-slots}   & patient, secretaire, admin & Créneaux libres d'un médecin. \\
  \method{GET}   & \texttt{/doctor/daily}      & medecin                    & Planning journalier du médecin. \\
  \method{PATCH} & \texttt{/:id/cancel}        & Authentifié                & Annule un rendez-vous. \\
  \method{PATCH} & \texttt{/:id/reschedule}    & Authentifié                & Reprogramme un rendez-vous. \\
  \method{POST}  & \texttt{/unavailability}    & medecin                    & Bloque un créneau. \\
  \method{PATCH} & \texttt{/:id/confirm}       & secretaire, admin          & Confirme un rendez-vous. \\
  \method{PATCH} & \texttt{/:id/no-show}       & secretaire, admin          & Marque un patient absent. \\
  \bottomrule
\end{tabular}
\end{center}

\section{Médecins — \texttt{/api/doctors}}

\begin{center}
\small
\begin{tabular}{@{} l p{4.5cm} p{4cm} p{4.5cm} @{}}
  \toprule
  \textbf{Méth.} & \textbf{Endpoint} & \textbf{Rôles} & \textbf{Description} \\
  \midrule
  \method{GET}    & \texttt{/search}       & Tous les rôles & Recherche par spécialité, langue, nom. \\
  \method{GET}    & \texttt{/me}           & medecin        & Profil du médecin connecté. \\
  \method{PATCH}  & \texttt{/me}           & medecin        & Mise à jour du profil. \\
  \method{POST}   & \texttt{/me/leaves}    & medecin        & Ajout d'un congé. \\
  \method{GET}    & \texttt{/me/leaves}    & medecin        & Liste des congés. \\
  \method{DELETE} & \texttt{/me/leaves/:lid} & medecin      & Suppression d'un congé. \\
  \bottomrule
\end{tabular}
\end{center}

\section{Dossiers Médicaux — \texttt{/api/records}}

\begin{center}
\small
\begin{tabular}{@{} l p{5.5cm} p{3.2cm} p{3.8cm} @{}}
  \toprule
  \textbf{Méth.} & \textbf{Endpoint} & \textbf{Rôles} & \textbf{Description} \\
  \midrule
  \method{POST}  & \texttt{/}                    & medecin              & Ajout d'une consultation. \\
  \method{PATCH} & \texttt{/my-record}           & patient              & Mise à jour antécédents / allergies. \\
  \method{POST}  & \texttt{/upload}              & patient              & Dépôt d'un document (20 Mo max). \\
  \method{POST}  & \texttt{/upload/:patientId}   & medecin              & Dépôt pour un patient tiers. \\
  \method{GET}   & \texttt{/my-record}           & patient              & Consultation du dossier personnel. \\
  \method{GET}   & \texttt{/my-record/}%
                   \texttt{consultations/:id/}%
                   \texttt{prescription/pdf}      & patient              & Ordonnance PDF. \\
  \method{GET}   & \texttt{/patient/:patientId}  & medecin, secret., admin & Dossier d'un patient tiers. \\
  \bottomrule
\end{tabular}
\end{center}

\section{Factures — \texttt{/api/invoices}}

\begin{center}
\small
\begin{tabular}{@{} l p{5cm} p{3.5cm} p{4cm} @{}}
  \toprule
  \textbf{Méth.} & \textbf{Endpoint} & \textbf{Rôles} & \textbf{Description} \\
  \midrule
  \method{POST}  & \texttt{/}                    & secretaire, admin & Création d'une facture. \\
  \method{GET}   & \texttt{/mine}                & patient           & Factures du patient connecté. \\
  \method{GET}   & \texttt{/}                    & secretaire, admin & Toutes les factures. \\
  \method{GET}   & \texttt{/overdue}             & secretaire, admin & Factures impayées > 30 jours. \\
  \method{GET}   & \texttt{/:id/pdf}             & patient, secret., admin & Téléchargement PDF. \\
  \method{GET}   & \texttt{/:id/feuille-soins}   & secretaire, admin & Feuille de soins CPAM. \\
  \method{POST}  & \texttt{/:id/send-email}      & secretaire, admin & Envoi par e-mail au patient. \\
  \method{PATCH} & \texttt{/:id/pay}             & secretaire, admin & Marque la facture payée. \\
  \bottomrule
\end{tabular}
\end{center}

\section{Administration — \texttt{/api/admin}}

\begin{center}
\small
\begin{tabular}{@{} l p{4.5cm} p{3.5cm} p{4.5cm} @{}}
  \toprule
  \textbf{Méth.} & \textbf{Endpoint} & \textbf{Rôles} & \textbf{Description} \\
  \midrule
  \method{POST}   & \texttt{/create-staff}  & admin             & Création d'un compte médecin ou secrétaire. \\
  \method{GET}    & \texttt{/staff}         & admin             & Liste du personnel. \\
  \method{PATCH}  & \texttt{/staff/:id}     & admin             & Mise à jour d'un membre du personnel. \\
  \method{DELETE} & \texttt{/accounts/:id}  & admin             & Suppression d'un compte. \\
  \method{GET}    & \texttt{/patients}      & secretaire, admin & Liste des patients. \\
  \method{GET}    & \texttt{/audit}         & admin             & Journal d'audit paginé. \\
  \bottomrule
\end{tabular}
\end{center}

\section{Établissement, Analytics et Avis}

\begin{center}
\small
\begin{tabular}{@{} l p{4.5cm} p{3.5cm} p{4.5cm} @{}}
  \toprule
  \textbf{Méth.} & \textbf{Endpoint} & \textbf{Rôles} & \textbf{Description} \\
  \midrule
  \method{GET}    & \texttt{/api/facility}          & Tous              & Informations de la clinique. \\
  \method{PUT}    & \texttt{/api/facility}          & admin             & Création ou mise à jour (upsert). \\
  \method{POST}   & \texttt{/api/facility/rooms}    & admin             & Ajout d'une salle. \\
  \method{DELETE} & \texttt{/api/facility/rooms/:rid}& admin            & Suppression d'une salle. \\
  \midrule
  \method{GET}    & \texttt{/api/analytics/kpis}    & admin             & KPIs : revenus, RDV, patients. \\
  \method{GET}    & \texttt{/api/analytics/revenue} & admin             & Revenus par médecin. \\
  \method{GET}    & \texttt{/api/analytics/export.pdf}  & admin         & Rapport PDF. \\
  \method{GET}    & \texttt{/api/analytics/export.xlsx} & admin         & Rapport XLSX. \\
  \midrule
  \method{POST}   & \texttt{/api/reviews}           & patient           & Dépôt d'un avis post-consultation. \\
  \method{GET}    & \texttt{/health}                & Public            & Vérification de santé du serveur. \\
  \bottomrule
\end{tabular}
\end{center}

% ============================================================
%  CHAPITRE 6 — SÉCURITÉ
% ============================================================
\chapter{Sécurité}

\section{Authentification et Autorisation}

\begin{itemize}[leftmargin=*,itemsep=3pt]
  \item \textbf{JWT HS256} — Tokens signés avec \texttt{JWT\_SECRET}. Durée de validité 7 jours. Le middleware \texttt{protect} vérifie la signature et charge l'utilisateur depuis MongoDB à chaque requête.
  \item \textbf{RBAC} — Le middleware \texttt{authorize(...roles)} compare \texttt{req.user.role} à la liste des rôles autorisés et retourne \texttt{403 Forbidden} en cas d'incompatibilité.
  \item \textbf{Hachage des mots de passe} — bcryptjs avec salt de 10 rondes, déclenché automatiquement par le hook \texttt{pre('save')} dès que le champ \texttt{password} est modifié.
  \item \textbf{TOTP 2FA} — Speakeasy génère un secret Base32. L'administrateur scanne un QR Code (qrcode), puis valide via un code 6 chiffres. Un \texttt{tempToken} JWT de courte durée est émis après le premier facteur ; le JWT final n'est délivré qu'après validation TOTP.
  \item \textbf{Google OAuth} — \texttt{google-auth-library} vérifie l'\texttt{idToken} côté serveur, évitant toute falsification côté client.
\end{itemize}

\section{Validation des Entrées}

\texttt{express-validator} valide les corps de requêtes. Le middleware \texttt{validate} centralise la réponse \texttt{422} avec le tableau d'erreurs. Les expressions régulières (format e-mail, complexité mot de passe, format date) sont définies directement dans les routes.

\section{Téléversement de Fichiers}

Multer est configuré avec un filtre MIME (PDF, JPEG, PNG, DICOM uniquement) et une limite de 20 Mo par fichier. Les fichiers sont stockés dans \texttt{uploads/} et servis en statique via Express.

\section{CORS}

L'origine autorisée est restreinte à \texttt{FRONTEND\_URL}. En son absence, la valeur par défaut \texttt{http://localhost:4200} est utilisée.

\section{Journal d'Audit (RGPD)}

Chaque accès sensible (connexion, consultation de dossier, création de compte, export) est enregistré dans la collection \texttt{AuditLog} via \texttt{auditLogger.js}. L'administrateur consulte ce journal paginé depuis l'interface d'administration.

% ============================================================
%  CHAPITRE 7 — INSTALLATION ET DÉMARRAGE
% ============================================================
\chapter{Installation et Démarrage}

\section{Prérequis}

\begin{itemize}[leftmargin=*,itemsep=2pt]
  \item Node.js \textbf{18 LTS} ou supérieur
  \item npm \textbf{9+}
  \item MongoDB \textbf{6+} (instance locale ou Atlas)
  \item Angular CLI : \texttt{npm install -g @angular/cli@17}
\end{itemize}

\section{Variables d'Environnement (Backend)}

Créer \texttt{backend/.env} à partir de \texttt{backend/.env.example} :

\begin{lstlisting}[style=bash]
PORT=3000
MONGODB_URI=mongodb://localhost:27017/medisync
JWT_SECRET=change_me_in_production_long_random_string
FRONTEND_URL=http://localhost:4200
# Optionnel - e-mails desactives si absent :
RESEND_API_KEY=re_xxxxxxxxxxxxxxxxxxxx
# Optionnel - Google OAuth desactive si absent :
GOOGLE_CLIENT_ID=xxxx.apps.googleusercontent.com
\end{lstlisting}

\begin{warnbox}[title=Important]
  Ne jamais commiter le fichier \texttt{.env}. Le \texttt{.gitignore} l'exclut par défaut.
\end{warnbox}

\section{Installation des Dépendances}

\begin{lstlisting}[style=bash]
# Backend
cd backend && npm install

# Frontend
cd ../frontend && npm install
\end{lstlisting}

\section{Démarrage des Services}

Ouvrir \textbf{trois terminaux distincts} :

\begin{lstlisting}[style=bash]
# Terminal 1 - Base de donnees MongoDB
mkdir -p ~/mongo-data
mongod --dbpath ~/mongo-data --port 27017

# Terminal 2 - Backend Express
cd backend
node seed.js     # Uniquement au premier lancement
node server.js   # ou "npm run dev" pour rechargement auto

# Terminal 3 - Frontend Angular
cd frontend
ng serve         # Ecoute sur http://localhost:4200
\end{lstlisting}

\begin{infobox}[title=Socket MongoDB bloqué]
  Si MongoDB ne démarre pas à cause d'un socket résiduel :\\
  \texttt{sudo rm /tmp/mongodb-27017.sock}
\end{infobox}

\section{Comptes de Démonstration}

Après \texttt{node seed.js}, les comptes suivants sont disponibles (mot de passe universel : \textbf{\texttt{Demo1234!}}) :

\begin{center}
\small
\begin{tabular}{@{} l l l l @{}}
  \toprule
  \textbf{Rôle} & \textbf{E-mail} & \textbf{Nom} & \textbf{Mot de passe} \\
  \midrule
  Administrateur & \texttt{admin@medisync.demo}              & Ali Benali          & \texttt{Demo1234!} \\
  Secrétaire     & \texttt{secretary@medisync.demo}          & Nadia Ouhbi         & \texttt{Demo1234!} \\
  \midrule
  Médecin 1      & \texttt{doctor@medisync.demo}             & Fatima Cherif       & \texttt{Demo1234!} \\
                 & \multicolumn{3}{l}{\footnotesize\textit{Cardiologie, Médecine générale — secteur 2}} \\[2pt]
  Médecin 2      & \texttt{doctor2@medisync.demo}            & Karim Amrani        & \texttt{Demo1234!} \\
                 & \multicolumn{3}{l}{\footnotesize\textit{Orthopédie, Traumatologie — secteur 2}} \\[2pt]
  Médecin 3      & \texttt{doctor3@medisync.demo}            & Leila Mansouri      & \texttt{Demo1234!} \\
                 & \multicolumn{3}{l}{\footnotesize\textit{Pédiatrie — secteur 1}} \\[2pt]
  Médecin 4      & \texttt{doctor4@medisync.demo}            & Omar Bensouda       & \texttt{Demo1234!} \\
                 & \multicolumn{3}{l}{\footnotesize\textit{Radiologie, Imagerie médicale — secteur 2}} \\[2pt]
  Médecin 5      & \texttt{doctor5@medisync.demo}            & Amina Hajji         & \texttt{Demo1234!} \\
                 & \multicolumn{3}{l}{\footnotesize\textit{Dermatologie, Vénérologie — secteur 2}} \\[2pt]
  \midrule
  Patient 1      & \texttt{marwane.elbaraka@gmail.com}       & Youssef Tazi        & \texttt{Demo1234!} \\
                 & \multicolumn{3}{l}{\footnotesize\textit{2 enfants : Nour, Hind}} \\[2pt]
  Patient 2      & \texttt{patient2@medisync.demo}           & Salma Idrissi       & \texttt{Demo1234!} \\
  Patient 3      & \texttt{patient3@medisync.demo}           & Hassan El Fassi     & \texttt{Demo1234!} \\
                 & \multicolumn{3}{l}{\footnotesize\textit{Conjoint : Khadija}} \\[2pt]
  Patient 4      & \texttt{patient4@medisync.demo}           & Zineb Moussaoui     & \texttt{Demo1234!} \\
  Patient 5      & \texttt{patient5@medisync.demo}           & Rachid Benkirane    & \texttt{Demo1234!} \\
                 & \multicolumn{3}{l}{\footnotesize\textit{Conjoint + 2 enfants : Amine, Rim}} \\
  \bottomrule
\end{tabular}
\end{center}

\section{Build de Production}

\begin{lstlisting}[style=bash]
cd frontend
ng build --configuration=production
# Les fichiers compiles se trouvent dans frontend/dist/
\end{lstlisting}

% ============================================================
%  CHAPITRE 8 — FONCTIONNALITÉS DÉTAILLÉES PAR RÔLE
% ============================================================
\chapter{Fonctionnalités Détaillées par Rôle}

\section{Portail Patient}

\begin{itemize}[leftmargin=*,itemsep=3pt]
  \item \textbf{Tableau de bord} : Statistiques personnelles (nombre de rendez-vous, prochaine consultation, factures en attente).
  \item \textbf{Recherche de médecins} : Filtrage par spécialité, langue et nom. Affichage des tarifs, secteur et disponibilités.
  \item \textbf{Prise de rendez-vous} : Sélection de médecin, calendrier, créneaux disponibles en temps réel, durée, motif.
  \item \textbf{Mes rendez-vous} : Historique paginé avec badges de statut, actions d'annulation.
  \item \textbf{Dossier médical} : Antécédents, allergies, consultations passées, ordonnances en PDF, pièces jointes.
  \item \textbf{Mes factures} : Historique des factures avec téléchargement PDF.
  \item \textbf{Avis} : Dépôt d'un avis noté (1–5 étoiles) après consultation terminée.
\end{itemize}

\section{Portail Médecin}

\begin{itemize}[leftmargin=*,itemsep=3pt]
  \item \textbf{Planning journalier} : Calendrier du jour avec code couleur par statut. Accès au formulaire de consultation directement depuis un RDV.
  \item \textbf{Formulaire de consultation} : Compte rendu structuré, prescriptions (médicament, posologie, durée), pièces jointes, changement de statut vers \textit{terminé}.
  \item \textbf{Congés} : Création de périodes de congés (dates + motif). Les congés bloquent automatiquement la prise de RDV pendant la période.
  \item \textbf{Profil} : Mise à jour des spécialités, langues, horaires, tarifs et actes CCAM.
\end{itemize}

\section{Portail Secrétaire}

\begin{itemize}[leftmargin=*,itemsep=3pt]
  \item \textbf{Planning global} : Vue tabulaire de l'ensemble des rendez-vous organisée par statut.
    Six onglets correspondent aux six états possibles d'un rendez-vous
    (\textit{en attente}, \textit{confirmé}, \textit{terminé}, \textit{annulé}, \textit{no-show},
    \textit{indisponible}). Dans chaque onglet, tous les rendez-vous du statut concerné sont
    affichés d'un seul tenant, quelle que soit leur date, et regroupés sous des en-têtes de date
    (ex. : \textit{Aujourd'hui}, \textit{Demain}, puis dates au format long en français) triés
    chronologiquement. Des badges de comptage sur les onglets \textit{en attente} et
    \textit{confirmé} signalent immédiatement les éléments nécessitant une action.
    Actions disponibles : confirmation, marquage no-show, report (rescheduling), annulation,
    émission de facture (onglet \textit{terminé}).
  \item \textbf{Patients} : Recherche, consultation et création de nouveaux dossiers.
  \item \textbf{Facturation} : Émission de factures pour les consultations terminées, suivi des impayés (> 30 j), envoi par e-mail, marquage de paiement.
\end{itemize}

\section{Portail Administrateur}

\begin{itemize}[leftmargin=*,itemsep=3pt]
  \item \textbf{Tableau de bord} : KPIs (revenus du mois, taux d'occupation, nouveaux patients). Raccourcis de navigation rapide.
  \item \textbf{Personnel} : CRUD des comptes médecins et secrétaires.
  \item \textbf{Établissement} : Informations de la clinique, gestion des salles et équipements.
  \item \textbf{Analytics} : Graphiques de revenus et répartition par médecin. Export PDF et XLSX.
  \item \textbf{Journal d'audit} : Consultation paginée des événements de sécurité tracés.
  \item \textbf{Sécurité 2FA} : Configuration du second facteur TOTP (QR Code + vérification).
\end{itemize}

% ============================================================
%  CHAPITRE 9 — MODULES FRONTEND
% ============================================================
\chapter{Modules et Services Frontend}

\section{Architecture Angular}

MediSync utilise l'architecture \textbf{Standalone Components} d'Angular 17 (sans NgModule). Le découpage est :

\begin{itemize}[leftmargin=*,itemsep=2pt]
  \item \textbf{Core} — Guards, intercepteurs, services, modèles, utilitaires. Aucune dépendance sur les features.
  \item \textbf{Features} — Cinq répertoires isolés : \texttt{auth}, \texttt{patient}, \texttt{practitioner}, \texttt{secretary}, \texttt{admin}. Chaque feature expose un composant shell avec lazy loading.
\end{itemize}

\section{Services Angular}

\begin{center}
\small
\begin{tabularx}{\textwidth}{@{} l X @{}}
  \toprule
  \textbf{Service} & \textbf{Responsabilités} \\
  \midrule
  \texttt{AuthService}        & Login, register, Google OAuth, 2FA, logout, \texttt{/auth/me}. \\
  \texttt{AppointmentService} & CRUD RDV, créneaux disponibles, planning journalier médecin. \\
  \texttt{DoctorService}      & Recherche médecins, profil propre, gestion des congés. \\
  \texttt{RecordService}      & Dossier médical, upload de fichiers, ordonnance PDF. \\
  \texttt{InvoiceService}     & Création, listage, statut, PDF, e-mail de factures. \\
  \texttt{AdminService}       & CRUD personnel, liste patients, journal d'audit. \\
  \texttt{FacilityService}    & Informations clinique, CRUD salles. \\
  \texttt{AnalyticsService}   & KPIs, revenus, export. \\
  \texttt{ReviewService}      & Dépôt d'avis post-consultation. \\
  \texttt{GoogleAuthService}  & Chargement script Google et popup OAuth. \\
  \bottomrule
\end{tabularx}
\end{center}

\section{Guards et Intercepteurs}

\begin{center}
\small
\begin{tabularx}{\textwidth}{@{} l X @{}}
  \toprule
  \textbf{Composant} & \textbf{Comportement} \\
  \midrule
  \texttt{AuthGuard}       & Redirige vers \texttt{/login} si aucun token JWT présent. \\
  \texttt{RoleGuard}       & Lit le rôle depuis le token décodé et redirige si incompatible avec la route. \\
  \texttt{AuthInterceptor} & Injecte \texttt{Authorization: Bearer <token>} sur chaque requête sortante. \\
  \texttt{ErrorInterceptor}& Intercepte les erreurs \texttt{401} et déclenche \texttt{AuthService.logout()}. \\
  \bottomrule
\end{tabularx}
\end{center}

% ============================================================
%  CHAPITRE 10 — SERVICES TRANSVERSAUX
% ============================================================
\chapter{Services Transversaux}

\section{Notifications E-mail}

Le service \texttt{emailService.js} utilise l'API Resend pour les e-mails transactionnels. Les templates HTML (confirmation, rappel 24h, facture) sont dans \texttt{emailTemplates.js}.

La fonction \texttt{sendNotification()} est appelée en mode \textit{fire-and-forget} (\texttt{.catch(() => \{\})}) pour éviter tout blocage des requêtes HTTP si le service e-mail est indisponible ou si la clé API est absente.

\section{Rappels Automatiques}

\texttt{reminderScheduler.js} utilise \texttt{node-cron} pour exécuter une tâche \textbf{toutes les heures} : elle interroge MongoDB pour trouver les rendez-vous confirmés dans les 24 heures à venir et envoie un e-mail de rappel au patient.

\section{Génération PDF}

\texttt{pdfGenerator.js} utilise PDFKit pour générer :
\begin{itemize}[leftmargin=*,itemsep=2pt]
  \item Les \textbf{factures} patient (entête clinique, détail acte, montant, statut).
  \item Les \textbf{feuilles de soins} (formulaire Sécurité Sociale).
  \item Les \textbf{ordonnances} avec prescriptions médicamenteuses.
  \item Les \textbf{rapports analytiques} pour l'administrateur.
\end{itemize}

\section{Export Excel}

\texttt{analyticsController.js} utilise ExcelJS pour générer des fichiers XLSX avec les données de revenus et KPIs, téléchargeables depuis l'interface d'administration.

% ============================================================
%  ANNEXES
% ============================================================
\appendix

\chapter{Codes de Statut HTTP}

\begin{center}
\small
\begin{tabular}{@{} r l p{7cm} @{}}
  \toprule
  \textbf{Code} & \textbf{Signification} & \textbf{Utilisation dans MediSync} \\
  \midrule
  200 & OK                    & Requête réussie \\
  201 & Created               & Ressource créée (RDV, facture, compte) \\
  400 & Bad Request           & Corps de requête invalide \\
  401 & Unauthorized          & Token absent, invalide ou expiré \\
  403 & Forbidden             & Rôle insuffisant \\
  404 & Not Found             & Ressource introuvable \\
  409 & Conflict              & Conflit (créneau déjà pris, doublon) \\
  422 & Unprocessable Entity  & Erreurs de validation (express-validator) \\
  500 & Internal Server Error & Erreur non gérée côté serveur \\
  \bottomrule
\end{tabular}
\end{center}

\chapter{Statuts des Rendez-vous}

\begin{center}
\small
\begin{tabular}{@{} l l p{7cm} @{}}
  \toprule
  \textbf{Statut} & \textbf{Acteur} & \textbf{Description} \\
  \midrule
  \texttt{en attente}  & Système              & État initial après création par un patient. \\
  \texttt{confirmé}    & Secrétaire / Admin   & RDV validé, rappel e-mail activé. \\
  \texttt{annulé}      & Tout acteur          & RDV annulé avant la consultation. \\
  \texttt{terminé}     & Médecin              & Consultation réalisée, dossier mis à jour. \\
  \texttt{no-show}     & Secrétaire / Admin   & Patient absent sans annulation. \\
  \texttt{indisponible}& Médecin              & Créneau bloqué (pas de patient associé). \\
  \bottomrule
\end{tabular}
\end{center}

\chapter{Structure du Token JWT}

Le payload JWT contient les informations minimales nécessaires à l'autorisation :

\begin{lstlisting}
{
  "id":   "ObjectId du compte Account",
  "role": "patient | medecin | secretaire | administrateur",
  "iat":  1716000000,
  "exp":  1716604800
}
\end{lstlisting}

La durée de validité est de \textbf{7 jours} (\texttt{7d}). Le secret de signature est \texttt{JWT\_SECRET} défini dans \texttt{.env}.

\end{document}
